\documentclass[12pt]{article}
\usepackage[margin=0.8in]{geometry}
\usepackage[T1]{fontenc}
\usepackage[utf8]{inputenc}
\usepackage{booktabs}
\usepackage{xcolor}
\usepackage{url}
\usepackage{array}
\usepackage{enumitem}
\usepackage{listings}

\definecolor{esfablue}{RGB}{28,68,135}
\definecolor{esfagreen}{RGB}{28,120,78}
\definecolor{lightgray}{RGB}{245,246,248}

\newcommand{\heading}[1]{\vspace{0.8em}\noindent{\Large\bfseries\color{esfablue}#1}\par\vspace{0.3em}}
\newcommand{\smallheading}[1]{\vspace{0.6em}\noindent{\bfseries\color{esfablue}#1}\par}
\lstset{breaklines=true,basicstyle=\ttfamily\small,columns=fullflexible}

\title{\textbf{ES-FA / BTP Presentation Report}\\
\large Event-Driven SNN FPGA Accelerator}
\author{Local smoke-run summary}
\date{\today}

\begin{document}
\maketitle

\begin{center}
\fcolorbox{esfablue}{lightgray}{
\begin{minipage}{0.92\linewidth}
\textbf{Current status:} software pipeline completed locally. Hardware/Vivado
execution was skipped in the latest smoke run. KV260 board measurement is the
main pending step.
\end{minipage}}
\end{center}

\heading{1. Project Idea}
ES-FA studies an event-driven FPGA execution path for spiking neural networks.
The basic idea is to use spike sparsity to reduce unnecessary memory access and
compute work. Instead of reporting only model accuracy, the project also tracks
hardware-facing metrics such as memory access, energy proxy, latency proxy, and
hardware validation outputs.

\heading{2. What Runs Now}
\begin{itemize}[leftmargin=*]
\item Baseline LIF SNN training runs locally.
\item Hardware-aware and adaptive experiments run locally.
\item Analysis generates ranking tables, evidence comparisons, and plots.
\item CLI demo works for control, math, and SNN-style classification queries.
\item Output sections are generated under \texttt{output/}.
\item Hardware RTL and validation scripts exist, but were skipped in this smoke run.
\end{itemize}

\heading{3. Current Architecture}
\begin{center}
\small
\begin{tabular}{p{0.25\linewidth}p{0.63\linewidth}}
\toprule
Layer & Current role \\
\midrule
\texttt{p1\_training/} & SNN model, training loop, baseline run, INT8/spike exports \\
\texttt{p2\_hardware\_model/} & spike, memory, energy, and latency proxy estimators \\
\texttt{experiments/} & hardware-aware loss and adaptive dataflow experiments \\
\texttt{analysis/} & ranking table, plots, evidence reports \\
\texttt{hardware/} & Verilog router, scheduler, event queue, memories, LIF PE \\
\texttt{hardware\_validation/} & xsim and Vivado flow for KV260-oriented validation \\
\texttt{deployment/} & CLI and adaptive execution policy \\
\texttt{output/} & report, paper, and presentation-ready artifacts \\
\bottomrule
\end{tabular}
\normalsize
\end{center}

\heading{4. Latest Smoke-Run Result}
\begin{center}
\small
\begin{tabular}{p{0.33\linewidth}rrrr}
\toprule
Run & Accuracy & Sparsity & Energy proxy & Score \\
\midrule
\texttt{exp1\_hardware\_aware\_loss} & 0.9570 & 0.5648 & 4,592,863.17 & 0.7699 \\
\texttt{baseline\_paper1} & 0.9570 & 0.5648 & 22,604,295.76 & 0.7387 \\
\texttt{exp5\_dataflow\_adaptation} & 0.9570 & 0.5648 & 45,016,894.73 & 0.6699 \\
\bottomrule
\end{tabular}
\normalsize
\end{center}

\smallheading{Main point}
The best current software run is \texttt{exp1\_hardware\_aware\_loss}. It keeps
95.70\% accuracy and gives a 79.68\% lower energy proxy compared with the dense
baseline estimate. This is a software proxy result, not measured board power.

\heading{5. Demo Commands}
\begin{lstlisting}
.\scripts\run_esfa_pipeline.ps1 -PythonExe .\.venv312\Scripts\python.exe `
  -Mode smoke -SkipVivado -SkipHardware

python deployment/cli_interface/main.py --query "hi"
python deployment/cli_interface/main.py --query "strike rate 167.55 balls 39"
python deployment/cli_interface/main.py --query "classify results/runtime/sample_signal.bin"
\end{lstlisting}

\heading{6. Hardware/Vivado Status}
The current RTL and validation stack are present, but the latest completed run
skipped hardware using \texttt{-SkipVivado -SkipHardware}. Existing local
hardware evidence rows are still available in the results folder, but the fresh
presentation result should be described as a software smoke run. Final hardware
claims should wait for:
\begin{itemize}[leftmargin=*]
\item xsim-only validation refresh;
\item full Vivado synthesis/implementation refresh;
\item physical KV260 board measurements.
\end{itemize}

\heading{7. What To Say In Presentation}
\begin{quote}
``The software and analysis side is running end-to-end locally. The best
hardware-aware run preserves 95.70\% validation accuracy and reduces the energy
proxy by 79.68\% compared with the dense baseline estimate. The RTL and
KV260-oriented validation scripts are in the repo, but fresh Vivado execution
and physical KV260 board measurements are the remaining hardware tasks.''
\end{quote}

\heading{8. Next Steps}
\begin{enumerate}[leftmargin=*]
\item Run xsim validation without the \texttt{-SkipHardware} flag.
\item Run Vivado validation when build time is available.
\item Boot KV260 and collect physical board measurements.
\item Compare board measurements with software estimators.
\item Tune estimator constants and rerun the report.
\end{enumerate}

\end{document}
